% GENERATED FILE. Edit canonical lessons, then run npm run generate:books.
% canonical-source-hash: fnv1a64:e267ad37a81138be
% canonical-lessons: PT-C446-arraial, PT-C446-romaria, PT-C446-procissao, PT-C446-padroeiro, PT-C446-desfilar

\chapter{Festas da terra}
\label{ch:pt-festas-da-terra}
\hlchaptermodality{446}

\begin{chapteropening}
\textbf{By the end of this chapter:} \emph{I can talk about a local festival: the open-air party, the pilgrimage festival, the procession, the patron saint and parading.}

Complete the last lesson of chapter 446: I can talk about a local festival: the open-air party, the pilgrimage festival, the procession, the patron saint and parading.
\end{chapteropening}

\section[o arraial]{o arraial — the open-air party (at a festival)}

\label{lesson:PT-C446-arraial}



\begin{quote}
Before the new one: say the Portuguese for to collapse, then the Portuguese for to knock down.
\end{quote}

\subsection*{You'll want to know: o arraial}
\textbf{o arraial} — \textquotedblleft{}the open-air party (at a festival)\textquotedblright{}.

\begin{quote}
No verão, há arraiais em todos os bairros.
\end{quote}

\begin{quote}
\emph{In summer, there are street parties in every neighbourhood.}
\end{quote}

In June the city quarters fill with \emph{arraiais}, music and grilled sardines.

\begin{cousinweb}
The origin of \emph{arraial} is not settled; it first meant a military camp and is usually linked with \emph{real} 'army camp', but EtymDB has no entry for it.
\end{cousinweb}

\begin{grammarlens}[title={What you've built}]
One more word for telling, explaining and describing at B1.
\end{grammarlens}

\subsection*{Your turn}
\begin{itemize}
\raggedright
  \item \emph{Say it:} \emph{o arraial}
  \item \emph{Say it:} \emph{o arraial}, once more
  \item \emph{Say it:} a sentence of your own with \emph{o arraial}
  \item \emph{Recall:} say the Portuguese for to collapse, then the Portuguese for to knock down, then say \emph{o arraial} again
\end{itemize}

\begin{tcolorbox}[breakable,colback=teal!4,colframe=teal!35!black,title={Before you move on}]
What does o arraial mean? (\textquotedblleft{}The open-air party (at a festival)\textquotedblright{}.) Say it once more.
\end{tcolorbox}
\section[a romaria]{a romaria — the pilgrimage festival, the pilgrimage}

\label{lesson:PT-C446-romaria}



\begin{quote}
Before the new one: say the Portuguese for to knock down, then the Portuguese for the open-air party.
\end{quote}

\subsection*{You'll want to know: a romaria}
\textbf{a romaria} — \textquotedblleft{}the pilgrimage festival, the pilgrimage\textquotedblright{}.

\begin{quote}
A aldeia inteira vai à romaria.
\end{quote}

\begin{quote}
\emph{The whole village goes to the pilgrimage festival.}
\end{quote}

A trip to a chapel, then food, music and dancing.

\begin{cousinweb}
\emph{romaria} is built on \emph{romeiro} 'pilgrim' with \emph{-ia} (EtymDB lists \emph{romaria} without giving a source).
\end{cousinweb}

\begin{grammarlens}[title={What you've built}]
One more word for telling, explaining and describing at B1.
\end{grammarlens}

\subsection*{Your turn}
\begin{itemize}
\raggedright
  \item \emph{Say it:} \emph{a romaria}
  \item \emph{Say it:} \emph{a romaria}, once more
  \item \emph{Say it:} a sentence of your own with \emph{a romaria}
  \item \emph{Recall:} say the Portuguese for to knock down, then the Portuguese for the open-air party, then say \emph{a romaria} again
\end{itemize}

\begin{tcolorbox}[breakable,colback=teal!4,colframe=teal!35!black,title={Before you move on}]
What does a romaria mean? (\textquotedblleft{}The pilgrimage festival, the pilgrimage\textquotedblright{}.) Say it once more.
\end{tcolorbox}
\section[a procissão]{a procissão — the procession}

\label{lesson:PT-C446-procissao}



\begin{quote}
Before the new one: say the Portuguese for the open-air party, then the Portuguese for the pilgrimage festival.
\end{quote}

\subsection*{You'll want to know: a procissão}
\textbf{a procissão} — \textquotedblleft{}the procession\textquotedblright{}.

\begin{quote}
A procissão passa pela rua principal.
\end{quote}

\begin{quote}
\emph{The procession goes along the main street.}
\end{quote}

\emph{A procissão ainda vai no adro} — it is early days yet.

\begin{cousinweb}
\emph{procissão} comes from Latin \emph{prōcessiō} 'marching forward, religious procession' (EtymDB), from \emph{prōcēdere} 'to go forward'.
\end{cousinweb}

\begin{grammarlens}[title={What you've built}]
One more word for telling, explaining and describing at B1.
\end{grammarlens}

\subsection*{Your turn}
\begin{itemize}
\raggedright
  \item \emph{Say it:} \emph{a procissão}
  \item \emph{Say it:} \emph{a procissão}, once more
  \item \emph{Say it:} a sentence of your own with \emph{a procissão}
  \item \emph{Recall:} say the Portuguese for the open-air party, then the Portuguese for the pilgrimage festival, then say \emph{a procissão} again
\end{itemize}

\begin{tcolorbox}[breakable,colback=teal!4,colframe=teal!35!black,title={Before you move on}]
What does a procissão mean? (\textquotedblleft{}The procession\textquotedblright{}.) Say it once more.
\end{tcolorbox}
\section[o padroeiro]{o padroeiro — the patron saint, the patron}

\label{lesson:PT-C446-padroeiro}



\begin{quote}
Before the new one: say the Portuguese for the pilgrimage festival, then the Portuguese for the procession.
\end{quote}

\subsection*{You'll want to know: o padroeiro}
\textbf{o padroeiro} — \textquotedblleft{}the patron saint, the patron\textquotedblright{}.

\begin{quote}
A festa do padroeiro é em agosto.
\end{quote}

\begin{quote}
\emph{The patron saint's festival is in August.}
\end{quote}

Feminine: \emph{a padroeira}.

\begin{cousinweb}
\emph{padroeiro} goes back to Latin \emph{patrōnus} 'protector' (EtymDB has \emph{patrono}), built on \emph{pater} 'father'.
\end{cousinweb}

\begin{grammarlens}[title={What you've built}]
One more word for telling, explaining and describing at B1.
\end{grammarlens}

\subsection*{Your turn}
\begin{itemize}
\raggedright
  \item \emph{Say it:} \emph{o padroeiro}
  \item \emph{Say it:} \emph{o padroeiro}, once more
  \item \emph{Say it:} a sentence of your own with \emph{o padroeiro}
  \item \emph{Recall:} say the Portuguese for the pilgrimage festival, then the Portuguese for the procession, then say \emph{o padroeiro} again
\end{itemize}

\begin{tcolorbox}[breakable,colback=teal!4,colframe=teal!35!black,title={Before you move on}]
What does o padroeiro mean? (\textquotedblleft{}The patron saint, the patron\textquotedblright{}.) Say it once more.
\end{tcolorbox}
\section[desfilar]{desfilar — to parade, to march past}

\label{lesson:PT-C446-desfilar}



\begin{quote}
Before the new one: say the Portuguese for the procession, then the Portuguese for the patron saint.
\end{quote}

\subsection*{You'll want to know: desfilar}
\textbf{desfilar} — \textquotedblleft{}to parade, to march past\textquotedblright{}.

\begin{quote}
As bandas desfilam pela avenida.
\end{quote}

\begin{quote}
\emph{The bands parade along the avenue.}
\end{quote}

Models \emph{desfilam} on a catwalk too.

\begin{cousinweb}
\emph{desfilar} was borrowed from French \emph{défiler} 'to march in file', built on \emph{fil} 'thread', from Latin \emph{fīlum} (Lewis \& Short).
\end{cousinweb}

\begin{grammarlens}[title={What you've built}]
One more word for telling, explaining and describing at B1.
\end{grammarlens}

\subsection*{Your turn}
\begin{itemize}
\raggedright
  \item \emph{Say it:} \emph{desfilar}
  \item \emph{Say it:} \emph{desfilar}, once more
  \item \emph{Say it:} a sentence of your own with \emph{desfilar}
  \item \emph{Recall:} say the Portuguese for the procession, then the Portuguese for the patron saint, then say \emph{desfilar} again
\end{itemize}

\begin{tcolorbox}[breakable,colback=teal!4,colframe=teal!35!black,title={Before you move on}]
What does desfilar mean? (\textquotedblleft{}To parade, to march past\textquotedblright{}.) Say it once more.
\end{tcolorbox}
